\section{Endpoint exponent audit and the revised loss ledger}
\label{sec:tpc59-endpoint}

The critical cutoff removes a geometrically forced positive-power
displacement, but it does not improve the inherited determinant-ladder
capacity or create singleton energy.  This section performs the complete
endpoint calculation for a general cutoff exponent and then replaces the
single relocation symbol of TPC--58 by the two logically distinct gates
identified above.

\subsection{The general \texorpdfstring{\(\xi\)}{xi} exponent dictionary}

Write
\begin{equation}
 \omega_X=\frac yQ=X^{\xi+o(1)},
 \qquad 0\le\xi<\frac{133}{400}.
 \label{eq:tpc59-endpoint-xi}
\end{equation}
The upper restriction merely keeps the compatible cofactor ceiling
polynomially growing; the identities below remain valid, with a bounded or
empty cofactor range, at and beyond that endpoint.  From
\eqref{eq:tpc59-endpoint-QJD} and
\cref{thm:tpc59-cutoff-cofactor-reciprocity}, the largest compatible
cofactor exponent is
\begin{equation}
 \theta_c(\xi)
 :=\frac{133}{400}-\xi.
 \label{eq:tpc59-critical-cofactor-exponent}
\end{equation}
The source-prime exponent and the high-cofactor singleton threshold
exponent are, respectively,
\begin{align}
 \lambda_L
 &:=\frac{267}{400}-\frac{10049}{52500}
   =\frac{99979}{210000},
 \label{eq:tpc59-endpoint-lambda-L}\\
 \lambda_{LD^2}
 &:=\lambda_L+2\frac{10049}{52500}
   =\frac{180371}{210000}.
 \label{eq:tpc59-endpoint-lambda-LD2}
\end{align}
For later reference, a compatible very-high residual label \(s=d_mr\)
obeys
\begin{equation}
 s\le X^{\lambda_D+\theta_c(\xi)+o(1)}
  =X^{110021/210000-\xi+o(1)},
 \qquad \lambda_D:=\frac{10049}{52500}.
 \label{eq:tpc59-compatible-residual-exponent}
\end{equation}
Thus every residual label remains of polynomial size, as required by the
TPC--55/57 divisor bounds \citep{WangTPC55,WangTPC57}.
Let \(M_X(R)\) denote the corresponding coefficient-blind maximum number
of terminal atoms in one residual fiber on a block \(r\asymp R\).

\begin{theorem}[Endpoint audit on every compatible cofactor block]
\label{thm:tpc59-compatible-endpoint-audit}
Let \(R=X^{\theta_R+o(1)}\) be a dyadic cofactor block contained in the
very-high compatible range, so that
\begin{equation}
 0\le\theta_R\le\theta_c(\xi).
 \label{eq:tpc59-compatible-theta-range}
\end{equation}
Then the following statements hold.
\begin{enumerate}[label=\textup{(\roman*)},leftmargin=3em]
 \item The source-prime/cofactor ratio satisfies
 \begin{equation}
  \frac RL
  \le X^{-15077/105000-\xi+o(1)}.
  \label{eq:tpc59-compatible-R-over-L}
 \end{equation}
 In particular \(R=o(L)\), uniformly over the compatible range.

 \item The exponent supplied by the coefficient-blind determinant-ladder
 capacity certificate is
 \begin{equation}
  \mu_{\rm cap}(R)
  :=(\lambda_L-\theta_R)_+
  =\lambda_L-\theta_R,
  \label{eq:tpc59-compatible-capacity-exponent}
 \end{equation}
 and therefore
 \begin{equation}
  \boxed{
  \mu_{\rm cap}(R)
  \ge\frac{15077}{105000}+\xi.}
  \label{eq:tpc59-compatible-capacity-floor}
 \end{equation}
 Its gap above the nominal endpoint allowance is
 \begin{equation}
  \boxed{
  \mu_{\rm cap}(R)-\frac1{400}
  \ge\frac{29629}{210000}+\xi.}
  \label{eq:tpc59-compatible-capacity-budget-gap}
 \end{equation}

 \item The TPC--57 same-divisor exclusion condition \(R>HL\), where
 \(H=\rad|h_0|=X^{o(1)}\), fails throughout the compatible range.  The
 corresponding large-exchanged-factor lower bound is therefore
 exponent-trivial there.

 \item The automatic full-fiber singleton condition
 \begin{equation}
  R>48HLD^2
  \label{eq:tpc59-automatic-singleton-condition}
 \end{equation}
 is separated from the largest compatible block by the fixed power
 \begin{equation}
  \boxed{
  \lambda_{LD^2}-\theta_c(\xi)
  =\frac{55273}{105000}+\xi.}
  \label{eq:tpc59-automatic-singleton-gap}
 \end{equation}
 Consequently the automatic-singleton range is empty on the present
 very-high carrier for every \(\xi\ge0\).
\end{enumerate}
\end{theorem}

\begin{proof}
The right side of \eqref{eq:tpc59-compatible-theta-range} and
\eqref{eq:tpc59-endpoint-lambda-L} give
\[
 \theta_R-\lambda_L
 \le\frac{133}{400}-\xi-\frac{99979}{210000}
 =-\frac{15077}{105000}-\xi,
\]
which proves \eqref{eq:tpc59-compatible-R-over-L}.  Since the right side is
strictly negative, \((\lambda_L-\theta_R)_+=\lambda_L-\theta_R\), and its
minimum occurs at \(\theta_R=\theta_c(\xi)\).  Hence
\[
 \lambda_L-\theta_R
 \ge\frac{99979}{210000}
       -\left(\frac{133}{400}-\xi\right)
 =\frac{15077}{105000}+\xi.
\]
Subtracting \(1/400=525/210000\) proves
\eqref{eq:tpc59-compatible-capacity-budget-gap}.  Because \(H=X^{o(1)}\),
\eqref{eq:tpc59-compatible-R-over-L} rules out \(R>HL\).
Finally,
\[
 \frac{180371}{210000}
 -\left(\frac{133}{400}-\xi\right)
 =\frac{55273}{105000}+\xi,
\]
which proves the last assertion.  Fixed constants such as \(48H\) are
absorbed by the displayed \(X^{o(1)}\) term.
\end{proof}

The capacity calculation is a statement about the strength of an
\emph{upper} certificate:
\begin{equation}
 M_X(R)\ll X^{o(1)}\left(1+\frac LR\right).
 \label{eq:tpc59-capacity-certificate-reminder}
\end{equation}
It does not assert that the actual multiplicity is as large as the right
side.  More importantly, even a degree-two carrier may have zero singleton
energy.  Thus \eqref{eq:tpc59-compatible-capacity-budget-gap} kills a
support-only transversal closure at the \(1/400\) endpoint; it is not a
value that may be entered as a proved singleton-occupancy loss.  A weighted
distinguished-vector estimate can still be strictly better, but none is
proved here.

\begin{corollary}[Critical cutoff is support-optimal]
\label{cor:tpc59-critical-cutoff-optimal}
Under the declared information \(\Delta_M=O(Q)\), and at the level of
fixed-power exponents, among the certified polynomial schedules
\(\omega_X=X^{\xi+o(1)}\), \(\xi\ge0\), the critical choice
\begin{equation}
 \boxed{\xi=0,\qquad y\asymp Q}
 \label{eq:tpc59-xi-zero-optimal}
\end{equation}
maximizes the compatible cofactor exponent and simultaneously minimizes
the certified capacity exponent, the failure exponent in \(R/L\), and the
gap to the automatic-singleton threshold.  Every fixed \(\xi>0\) worsens
all four quantities by exactly \(\xi\).
\end{corollary}

This is an endpoint optimization only.  It proves neither that a critical
block is present in every active row nor that any such block terminalizes.
Sharper support information such as \(\Delta_M=o(Q)\) could permit a
different sub-\(Q\) cutoff analysis; no such information is assumed here.
Those questions are measured exactly by
\(\beta_{\cR}\), \(\rho_{\cR}(\zeta)\), and \(\tau_{\cR}\) in
\cref{sec:tpc59-occupancy}.

\subsection{Retiring relocation: feasibility before occupancy}

TPC--58 used one symbol \(\lambda_{\rm reloc}\) for the passage from a
parent carrier to the smaller cofactor range produced by its fixed-
\(\eta\) cutoff \citep{WangTPC58}.  That notation mixed a support
feasibility question with an energy-transfer estimate.  They are not
additive loss exponents.  The correct replacement is the ordered pair
\begin{equation}
 \boxed{
 \lambda_{\rm reloc}
 \quad\rightsquigarrow\quad
 \{\text{cutoff feasibility prerequisite}\}
 \ \text{followed by}\ \lambda_{\rm occ}.}
 \label{eq:tpc59-relocation-split}
\end{equation}
The feasibility prerequisite is the trichotomy in
\cref{thm:tpc59-cutoff-cofactor-reciprocity}: a candidate block is safe,
lies in the bounded boundary region, or is empty.  It has no associated
lower-energy inequality.  At the critical cutoff, blocks at exponent
\(133/400\) are feasible, so there is no separate geometric exponent to
enter in an energy ledger.  For the fixed-\(\eta\) cutoff, by contrast, the
old \(R\asymp J\) block is eventually empty; this is a closed gate, not an
energy loss equal to \(\eta\).

Once feasibility is established, the only parent-to-window energy charge
is the actual occupancy exponent.  A typical distinguished input would be
\begin{equation}
 \cE_{\cR}(\zeta)
 \ge X^{-\lambda_{\rm occ}+o(1)}\cE_{\rm par}(\zeta),
 \label{eq:tpc59-occupancy-loss-input}
\end{equation}
or the corresponding coefficient-uniform assertion
\(\beta_{\cR}\ge X^{-\lambda_{\rm occ}+o(1)}\).  Neither follows from
cutoff admissibility.

There is no hidden transport map behind
\eqref{eq:tpc59-occupancy-loss-input}.  It is an orthogonal coordinate
projection onto cofactor labels already present in the parent carrier.
Moreover, factor exchange inside one residual fiber satisfies
\(r'/r\in[1/2,2]\).  It cannot move energy across the fixed-power gap
created by a cutoff with \(\xi>0\).

\subsection{The complete prospective endpoint ledger}

Let \(\cD_{\rm up}\) denote an explicitly declared upstream raw atomic
energy.  If a packet-selection theorem is used, its normalization is
\begin{equation}
 \cE_{\rm par}(\zeta)
 \ge X^{-\lambda_{\rm sel}+o(1)}\cD_{\rm up},
 \label{eq:tpc59-upstream-selection-normalization}
\end{equation}
where \(\cE_{\rm par}\) is exactly the one-side slice energy defined in
\eqref{eq:tpc59-parent-energy}.  For one distinguished coefficient vector,
write the remaining conditional chain before inserting its exponents as
\begin{equation}
 \begin{aligned}
 \cD_{\rm up}
 &\xrightarrow{\rm selection}\cE_{\rm par}
 \xrightarrow{\rm occupancy}\cE_{\cR}\\
 &\xrightarrow{\rm terminal}\cD_{\cR}^{\rm term}
 \hookrightarrow\cD_H
 \xrightarrow{\rm dyadic}\cD_{k_*}\\
 &\xrightarrow{\rm singleton}\cG_H(y;\zeta)
 \xrightarrow{\rm native\ erasure}\mathcal N_H
 \xrightarrow{\rm affine}\mathcal N_{\rm aff}^{\rm rep}\\
 &
 \xrightarrow{\rm off\mbox{-}carrier\ reassembly}
 \mathcal N_{\rm slice}^{\rm full}\\
 &
 \xrightarrow{\rm boundary/tail}\text{parent arithmetic target}.
 \end{aligned}
 \label{eq:tpc59-prospective-transfer-chain}
\end{equation}
Every arrow denotes a different assertion.  The unlabeled inclusion is
exact because the indicator defining \(\cD_{\cR}^{\rm term}\) already
imposes \(p>y\).  Terminal activation includes primality, the shaped
terminal cutoff, squarefree and coprimality conditions; it is not part of
the cofactor projection.  The high-only square \(\cG_H\), not the terminal
square containing \(p\le y\), is the input to the TPC--58 native-erasure
map.

Here the TPC--58 maps \(T_H\) and \(C_{\mu,H}\) are understood on the same
represented terminal/Walsh baseline \(\Ran\iota\) used to define
\(\cG_H\).  Thus \(\mathcal N_{\rm aff}^{\rm rep}\) contains the base
column and the represented high column, but not the missing-high coordinates
measured by \(\cD_{\rm miss}\) in
\eqref{eq:tpc59-missing-high-Walsh-energy}.  Restoring those coordinates is
the displayed off-carrier reassembly arrow.  It is a genuine lower-transfer
gate: an omitted high vector can cancel the represented affine vector after
restoration, so no norm monotonicity is available.
If a future lower transfer proves
\begin{equation}
 \mathcal N_{\rm slice}^{\rm full}
 \ge X^{-\lambda_{\rm reasm}+o(1)}
       \mathcal N_{\rm aff}^{\rm rep},
 \label{eq:tpc59-reassembly-loss-input}
\end{equation}
then \(\lambda_{\rm reasm}\) is its separate ledger charge.

The least maximizing dyadic index \(k_*=k_*(\zeta)\) is allowed to depend
on this distinguished vector.  Thus the zero block exponent in this chain
is not a coefficient-uniform theorem for one fixed dyadic cell.  Such a
uniform route would have to fix the cell independently of \(\zeta\) or
prove a separate finite-cover lower reassembly.

If future work proves polynomial lower bounds for all arrows in
\eqref{eq:tpc59-prospective-transfer-chain}, the complete exponent is
\begin{equation}
 \boxed{
 \begin{aligned}
 \Lambda_{\rm tot}:={}&
 \lambda_{\rm sel}+\lambda_{\rm occ}
 +\lambda_{\rm blk}
 +\lambda_{\rm term}+\lambda_{\rm sing}\\
 &+\lambda_H+\lambda_{\rm aff}
 +\lambda_{\rm reasm}
 +\lambda_{\rm bdry}+\lambda_{\rm tail}.
 \end{aligned}}
 \label{eq:tpc59-complete-loss-ledger}
\end{equation}
The critical cutoff establishes feasibility at exponent \(R\asymp J\)
without creating an energy-ledger entry.  The one exact zero fixed-power
energy-localization cost established here is
\begin{equation}
 \boxed{\lambda_{\rm blk}=0,}
 \label{eq:tpc59-zero-block-loss}
\end{equation}
as a distinguished-vector statement relative to an already activated
nonzero very-high terminal diagonal, as in
\cref{thm:tpc59-exact-dyadic-localization}.  No positive lower scale or
numerical loss exponent is proved for the remaining arithmetic gates.

The safe endpoint requirement is the strict inequality
\begin{equation}
 \boxed{\Lambda_{\rm tot}<\frac1{400}.}
 \label{eq:tpc59-strict-endpoint-ledger}
\end{equation}
Equality is not accepted without a constant-level comparison and quantified
remainders.  In particular, the logarithmic factor in the dyadic
pigeonhole is \(X^{o(1)}\) and has zero fixed exponent, but it can still
spoil a closure attempted exactly at \(1/400\).

\subsection{Level and kill-test boundary}

The conclusions of this section have the following precise status.
\begin{description}[leftmargin=3.3em,style=nextline]
 \item[Exact L0/L1 progress.]
 The rational endpoint identities, cutoff optimization, support-capacity
 audit, empty automatic-singleton range, feasibility/occupancy split, and
 strict loss ledger are exact.  Their use of \(s=d_mr\), the literal fixed
 shift, and raw cofactor atoms is the L1 specialization.

 \item[Closed support-only doors.]
 On every compatible block, the available transversal/capacity exponent
 already exceeds \(1/400\), while the deterministic same-divisor and
 automatic-singleton thresholds are outside the carrier.  Increasing
 \(\xi\) cannot repair either failure.

 \item[Open arithmetic doors.]
 A positive distinguished or uniform cofactor occupancy, terminal
 activation, weighted singleton fraction, native-erasure angle,
 represented affine noncancellation angle, missing-high/off-carrier
 reassembly, and boundary/tail reconnection remain independent hypotheses.

 \item[Not proved (L2).]
 No positive fixed-\(h_0\) arithmetic lower bound, parity improvement,
 prime-pair asymptotic, or twin-prime conclusion follows from the endpoint
 optimization.
\end{description}

Thus the critical cutoff genuinely removes one artificial positive-power
support displacement and the dyadic theorem removes a second bookkeeping
loss.  The remaining obstruction is no longer cutoff geometry: it is the
 actual coefficient-weighted occupancy and activation of the compatible
 carrier, followed by the two undeleted grouping gates and restoration of
 the unrepresented high coordinates.
